% ====================================================================
% DATEI: content/09_thresholds.tex
% Erweiterte Schwellenwertlandschaft — Core 1.3.3 (Abschlussversion)
% ====================================================================

\section{Erweiterte Schwellenwertlandschaft}
\label{sec:thresholds-extended}

Dieses Kapitel rekonstruiert die vollständige Theorie der Schwellenwerte in der Ontologie der
Kontinua (OC) und konsolidiert strukturelle Elemente aus den Läufen Physik,
Chemie, Biologie und Kognition.
Schwellenwerte sind der zentrale Mechanismus, über den
qualitative Veränderungen, Dimensionalemergenz, Kollaps und Tod erzeugt werden.

Schwellenwerte sind universell: Sie gelten für physische, chemische, biologische,
kognitive, soziale und theoretische Kontinua im gesamten Hierarchiebereich
\(K_0\)–\(K_{10}\).

\subsection{Definition eines Schwellenwerts}

Ein \emph{Schwellenwert} wird durch eine Nebenbedingungsfunktion dargestellt durch
\[
   f_k : \Omega(K) \rightarrow \mathbb{R}
\]
mit der Eigenschaft
\[
   f_k(s) \le 0
   \quad\text{für alle zulässigen Zustände } s\in\Omega(K).
\]

Schwellenwerte trennen qualitativ unterschiedliche dynamische Regime:
\begin{itemize}
    \item \( f_k(s) < 0 \): sicheres Regime (Inneres),
    \item \( f_k(s) = 0 \): Randregime,
    \item \( f_k(s) > 0 \): verbotendes Regime (außerhalb von \(\Omega(K)\)).
\end{itemize}

Damit legen Schwellenwerte gemeinsam fest:
\begin{itemize}
    \item die Gestalt der zulässigen Menge \(\Omega(K)\),
    \item den Rand \(\partial\Omega(K)\),
    \item die Bedingungen, unter denen das Kontinuum existiert oder kollabiert.
\end{itemize}

Schwellenwerte sind keine statischen Konstanten.
Sie hängen von Potentialen \(P(t)\), Flüssen \(J(t)\), Achsen \(A(K)\), Umwelt-\
bedingungen und der Struktur des Einbettungsraums \(M\) ab.
Auf struktureller Ebene ist dies durch den Schwellenwert-Evolutionsoperator
\(H\) der allgemeinen Evolutionsmaschine kodiert:
\[
  \frac{d\Theta_k}{dt}
   = H_k\big(\Theta_k, P, J, A, M\big),
\]
wobei \(\Theta_k\) die Parametermenge bezeichnet, die die Nebenbedingung \(f_k\) bestimmt.

\subsection{Taxonomie der Schwellenwerte}

OC unterscheidet sieben universelle Klassen von Schwellenwerten. Diese Taxonomie wurde
in archival source corpus stabilisiert und in den Domänenläufen präzisiert.

\paragraph{1. Existenz-Schwellenwerte \(\Theta_{\rm exist}\).}
Bedingungen, die für \(\Omega(K)\neq\emptyset\) erforderlich sind.
Beispiele:
\begin{itemize}
    \item minimale Energiebedingungen für gebundene physische Zustände (\(K_2\)),
    \item minimale katalytische Verschlussbedingung (RAF-Existenz) in \(K_3\),
    \item minimale Membrankontinuität und Integrität in \(K_4\),
    \item minimale Kohärenz interner Modelle in \(K_6\),
    \item minimale Legitimität und Vertrauen für Institutionen in \(K_7\).
\end{itemize}

\paragraph{2. Stabilitäts-Schwellenwerte \(\Theta_{\rm stab}\).}
Bedingungen, unter denen Flüsse beschränkt bleiben und Zyklen stabil sind.
Beispiele:
\begin{itemize}
    \item Konfinierung und Beschränktheitsbedingungen in der Physik (\(K_2\)),
    \item osmotische Homöostasebedingungen in Protokellen (\(K_4\)),
    \item Stabilität des Membranpotentials in Neuronen (\(K_5\)),
    \item normative Stabilität in sozialen Kontinua (\(K_7\)).
\end{itemize}

\paragraph{3. Kritische Schwellenwerte \(\Theta_{\rm crit}\).}
Oberflächen qualitativer Veränderungen bei fester Dimensionalität.
Beispiele:
\begin{itemize}
    \item Phasenübergänge und BKT-artige Schwellenwerte (\(K_2\)),
    \item Auftreten katalytischer Netzwerke (\(K_3\)),
    \item metabolischer Umschaltvorgang in \(K_4\),
    \item Initiierungsschwelle von Aktionsspitzen in \(K_5\),
    \item Vorhersagefehler-Bifurkationen in \(K_6\),
    \item institutionelle Bifurkationen in \(K_7\).
\end{itemize}

\paragraph{4. Dimensionalitäts-Schwellenwerte \(\Theta_{\rm dim}\).}
Nebenbedingungen, die die Erzeugung neuer Achsen auslösen.
Dimensionalitäts-Schwellenwerte sind zentral in OC:
\[
   T(K,t) > \Theta_{\rm dim}(K)
   \quad\Rightarrow\quad
   \text{Auftreten einer neuen Achse } A_{\rm new}.
\]

Sie steuern die Übergänge \(K_x \to K_{x+1}\) in der gesamten Hierarchie.

\paragraph{5. Todes-Schwellenwerte \(\Theta_{\rm death}\).}
Werte, jenseits derer kein zulässiger Zustand mehr existiert.
Äquivalent dazu gilt:
\[
   \forall s\in\overline{\Omega(K)}:\ \exists k\ \text{mit } f_k(s) > 0
   \quad\Longleftrightarrow\quad
   \Omega(K)=\emptyset.
\]

\paragraph{6. Expressivitäts-Schwellenwerte \(\Theta_{\rm expr}\).}
Schwellenwerte für Darstellungskapazität. Ein Kontinuum kollabiert, wenn
\[
   \dim\big(A(K)\big)
   <
   \dim\big({\rm Differences}(K)\big),
\]
wohingegen \(\mathrm{Differences}(K)\) die effektive Dimension strukturell relevanter
Unterscheidungen bezeichnet, die durch den Einbettungsraum und
Randbedingungen erzwungen werden. Diese Klasse ist für kognitive, soziale und
meta-theoretische Kollapse zentral.

\paragraph{7. Einbettungsschwellenwerte \(\Theta_{\rm embed}\).}
Bedingungen durch den Einbettungsraum \(M_x\).
Ein Kontinuum kann nicht als lebendes Kontinuum existieren, wenn der Einbettungsraum
seine Achsen, Schwellenwerte oder Flüsse nicht trägt; formal wird dies durch die
Kompatibilitätsbedingungen in Satz~4 und Satz~5 (Ergebniskapitel) erfasst.

\subsection{Schwellenwert-Geometrie}

Schwellenwerte induzieren eine geschichtete Geometrie auf \(\Omega(K)\).
Für jede Schwellenwertfunktion \(f_k\) betrachten wir die Niveaumengen
\[
  \Sigma_k(c) = \{\, s \in \Omega(K) \mid f_k(s) = c \,\}.
\]

Insbesondere:
\begin{itemize}
    \item \(\Sigma_k(0)\) definiert Komponenten des Randes \(\partial\Omega(K)\),
    \item \(\Sigma_k(c<0)\) definiert sichere Regionen,
    \item \(\Sigma_k(c>0)\) definiert verbotene Regionen.
\end{itemize}

Interaktionen zwischen Schwellenwerten erzeugen:
\begin{itemize}
    \item Mehrfach-Schwellenwert-Schnittmengen,
    \item Spitzen- und Faltstrukturen (katastrophentheoretische Geometrien),
    \item Regionen mit konzentrierter Krümmung,
    \item heterogene Struktur auf Patches am Rand
          (vgl. the boundary geometry chapter).
\end{itemize}

Die gemeinsame Schwellenwert-Geometrie bestimmt:
\begin{itemize}
    \item die globale Gestalt von \(\Omega(K)\),
    \item zulässige Trajektorien und Zyklusstruktur,
    \item die verfügbaren Übergänge des Kontinuums.
\end{itemize}

\subsection{Dimensionalitäts-Schwellenwerte}

Dimensionalitäts-Schwellenwerte nehmen in OC eine privilegierte Rolle ein, da sie
entscheiden, wann neue Unterschiede innerhalb der aktuellen Achsen nicht mehr
repräsentierbar sind.

Sei \(T(K,t)\) die strukturelle Spannung.
Dann gilt schematisch:
\[
  T(K,t) < \Theta_{\rm dim}(K)
   \quad\Rightarrow\quad
   \text{alle relevanten Unterschiede können innerhalb von } A(K) \text{ ausgedrückt werden},
\]
\[
  T(K,t) = \Theta_{\rm dim}(K)
   \quad\Rightarrow\quad
   \text{Projektion auf } A(K)\ \text{scheitert},
\]
\[
  T(K,t) > \Theta_{\rm dim}(K)
   \quad\Rightarrow\quad
   \text{eine neue Achse } A_{\rm new} \text{ ist strukturell erforderlich}.
\]

Dimensionalitäts-Schwellenwerte regeln insbesondere:
\begin{itemize}
    \item Entstehung von \(K_1\) aus der strukturellen Grundlage \(K_0\),
    \item Auftreten chemischer Achsen in \(K_3\),
    \item Membranachsen in \(K_4\),
    \item Erregbarkeitsachsen in \(K_5\),
    \item kognitive repräsentationale Achsen in \(K_6\),
    \item soziale normative Achsen in \(K_7\),
    \item meta-logische Achsen in \(K_9\)–\(K_{10}\).
\end{itemize}

\subsection{Schwellenwert-Kaskaden}

Schwellenwerte interagieren häufig in Kaskaden der Form
\[
  \Theta_{\rm grad}
  \;\longrightarrow\;
  \Theta_{\rm perm}
  \;\longrightarrow\;
  \Theta_{\rm mem},
\]
oder allgemeiner: Gradienten \(\to\) Permeabilität \(\to\) strukturelles Versagen.

Repräsentative Beispiele:

\begin{itemize}
    \item \textbf{Vesikelriss (\(K_4\)).}
          Osmotische Gradienten überschreiten \(\Theta_{\rm grad}\), die Permeabilität steigt
          (\(\Theta_{\rm perm}\)), und die Membranintegrität geht verloren
          (\(\Theta_{\rm mem}\)).
    \item \textbf{Aktionspotenziale (\(K_5\)).}
          Die Depolarisation überschreitet eine Gradienten-Schwelle, Kanäle öffnen sich bei
          \(\Theta_{\rm crit}\), und ein Spike-Zyklus wird aktiviert.
    \item \textbf{Institutioneller Kollaps (\(K_7\)).}
          Vertrauensdefizite überschreiten \(\Theta_{\rm grad}\), normative Kohärenz
          versagt (\(\Theta_{\rm expr}\)), und der institutionelle Tod folgt
          (\(\Theta_{\rm death}\)).
\end{itemize}

Solche Kaskaden sind zentral für frühe Lebensformen, neuronale Systeme sowie
soziale und infrastrukturelle Dynamiken.

\subsection{Todes-Schwellenwerte}

Todes-Schwellenwerte bestimmen, wann ein Kontinuum als lebendes Kontinuum aufhört zu
gelten.
Die folgenden strukturellen Bedingungen sind äquivalente Manifestationen von Tod:
\begin{itemize}
    \item vollständige Verletzung von Schwellenwerten in allen Zuständen:
          \(\forall s:\ \exists k \text{ mit } f_k(s) > 0\);
    \item Verschwinden unterstützender Zyklen: \(C(K)=\emptyset\);
    \item Kollaps des zulässigen Zustandsraums: \(\Omega(K)=\emptyset\);
    \item Repräsentationsversagen, z.\,B. \(\Theta_{\rm expr}\) überschritten;
    \item Unvereinbarkeit mit dem Einbettungsraum (keine Konfiguration erfüllt
          sowohl interne als auch Einbettungs-Bedingungen).
\end{itemize}

Sobald \(\Omega(K)=\emptyset\), ist der Tod strukturell irreversibel:
Kein Operator innerhalb derselben Ebene kann \(\Omega(K)\) erneut erzeugen;
jede scheinbare Erholung entspricht der Geburt eines neuen Kontinuums \(K'\).

\subsection{Beispiele über Ebenen}

Die Schwellenwert-Taxonomie manifestiert sich auf jeder Ebene der Hierarchie unterschiedlich.

\paragraph{\texorpdfstring{\(K_1\)}{K_1} (Geometrische Kontinua).}
\begin{itemize}
    \item \(\Theta_{\rm exist}\): Differenzierbarkeits- und Regularitätsbedingungen,
          die die klassische Region \(\Omega_{\rm cl}\) festlegen,
    \item \(\Theta_{\rm crit}\): klassische Bifurkationen und Stabilitätsverlust.
\end{itemize}

\paragraph{\texorpdfstring{\(K_2\)}{K_2} (Physikalische Kontinua).}
\begin{itemize}
    \item \(\Theta_{\rm crit}\): Phasenübergänge, einschließlich BKT-Schwellen,
    \item \(\Theta_{\rm dim}\): Kohärenz-Schwellenwerte im Zusammenhang mit
          Massenentstehung und Ordnung,
    \item \(\Theta_{\rm death}\): Konfinenz-/Entkopplungsgrenzen, an denen das
          betreffende Kontinuum aufhört zu existieren.
\end{itemize}

\paragraph{\texorpdfstring{\(K_3\)}{K_3} (Chemische Kontinua).}
\begin{itemize}
    \item \(\Theta_{\rm exist}\): katalytischer Abschluss und RAF-Existenz,
    \item \(\Theta_{\rm crit}\): Perkolation von Reaktionsnetzwerken,
    \item \(\Theta_{\rm expr}\): unzureichende chemische Achsen zur Repräsentation der erforderlichen
          Reaktionsvielfalt.
\end{itemize}

\paragraph{\texorpdfstring{\(K_4\)}{K_4} (Präbiotische/biologische Membrankontinua).}
\begin{itemize}
    \item \(\Theta_{\rm grad}\): osmotische Druckgrenzen,
    \item \(\Theta_{\rm perm}\): Permeabilitäts-Schwellenwerte,
    \item \(\Theta_{\rm mem}\): Membranruptur- und Krümmungsschwellen.
\end{itemize}

\paragraph{\texorpdfstring{\(K_5\)}{K_5} (Erregbare Kontinua).}
\begin{itemize}
    \item \(\Theta_{\rm crit}\): Auslöseschwellen für Spikes,
    \item \(\Theta_{\rm stab}\): Refraktärbedingungen und Erregbarkeitsgrenzen,
    \item \(\Theta_{\rm death}\): Verlust der Erregbarkeit und Zusammenbruch erregbarer Zyklen.
\end{itemize}

\paragraph{\texorpdfstring{\(K_6\)}{K_6} (Kognitive Kontinua).}
\begin{itemize}
    \item \(\Theta_{\rm expr}\): Grenzen der Bindungskapazität und repräsentationalen
          Komplexität,
    \item \(\Theta_{\rm crit}\): Schwellen für Vorhersagefehler,
    \item \(\Theta_{\rm stab}\): Gedächtniskohärenz und Modellstabilitätsgrenzen.
\end{itemize}

\paragraph{\texorpdfstring{\(K_7\)}{K_7}–\texorpdfstring{\(K_8\)}{K_8} (Soziale und zivilisatorische Kontinua).}
\begin{itemize}
    \item \(\Theta_{\rm exist}\): Schwellen für Vertrauen und Legitimität,
    \item \(\Theta_{\rm expr}\): institutionelle Kapazität und expressive Grenzen,
    \item \(\Theta_{\rm death}\): systemischer Kollaps und Zusammenbruch institutioneller
          sowie infrastruktureller Zyklen.
\end{itemize}

\paragraph{\texorpdfstring{\(K_9\)}{K_9}–\(K_{10}\) (Theoretische und meta-theoretische Kontinua).}
\begin{itemize}
    \item \(\Theta_{\rm expr}\): expressive Kapazität von Theorien und
          Metasprachen,
    \item \(\Theta_{\rm dim}\): Schwellen für Metaebenen-Expansion,
    \item \(\Theta_{\rm death}\): Inkonsistenz, Inkohärenz oder strukturelle
          Unvereinbarkeit im Modellraum.
\end{itemize}

\subsection{Zusammenfassung}

Die erweiterte Schwellenwertlandschaft liefert das strukturelle Rückgrat für alle
qualitativen Änderungen in OC. Schwellenwerte:
\begin{itemize}
    \item prägen \(\Omega(K)\) und definieren \(\partial\Omega(K)\),
    \item steuern Phasenübergänge, Emergenz neuer Dimensionen und Kollaps,
    \item vermitteln Kaskaden, die Gradienten, Permeabilität und strukturelles Versagen verbinden,
    \item kodieren expressive und Einbettungsbeschränkungen, die Kontinua töten können,
          selbst ohne energetische Überlastung.
\end{itemize}

Zusammen mit der Rand- und Patch-Geometrie aus
Abschnitt~the boundary geometry chapter vervollständigt die Schwellenwertlandschaft die
Beschreibung, wie Kontinua im OC-Rahmen entstehen, aufrechterhalten werden und
zerstört werden.
